\documentclass[11pt,a4paper]{article}

\usepackage[T1]{fontenc}
\usepackage[utf8]{inputenc}
\usepackage[spanish,es-noshorthands,es-tabla]{babel}
\usepackage{lmodern}
\usepackage{microtype}
\usepackage[margin=2.3cm]{geometry}
\usepackage[table]{xcolor}
\usepackage{booktabs}
\usepackage{tabularx}
\usepackage{array}
\usepackage{multirow}
\usepackage{enumitem}
\usepackage{tcolorbox}
\usepackage{titlesec}
\usepackage{amsmath,amssymb}
\usepackage{graphicx}
\usepackage{caption}
\usepackage{tikz}
\usepackage{pgfplots}
\pgfplotsset{compat=1.16}
\usepackage{tocloft}
\setlength{\cftsecnumwidth}{2em}
\renewcommand{\cftsecleader}{\cftdotfill{\cftdotsep}}
\usepackage[hidelinks]{hyperref}

\usetikzlibrary{babel,arrows.meta,positioning,calc,shapes.geometric}
\tcbuselibrary{skins,breakable}

\definecolor{azulnuevo}{HTML}{1F4E79}
\definecolor{grisrep}{HTML}{6B7280}
\definecolor{ambar}{HTML}{B45309}
\definecolor{verde}{HTML}{15803D}
\definecolor{fondonuevo}{HTML}{F3F7FB}
\definecolor{fondorep}{HTML}{F4F4F5}
\definecolor{fondoaviso}{HTML}{FFF7ED}
\definecolor{fondoej}{HTML}{F0FDF4}
\definecolor{critrojo}{HTML}{DC2626}
\definecolor{azultemp}{HTML}{1D4ED8}
\definecolor{rojotard}{HTML}{B91C1C}

\titleformat{\section}{\Large\bfseries\color{azulnuevo}}{\thesection}{0.6em}{}
\titleformat{\subsection}{\large\bfseries}{\thesubsection}{0.6em}{}
\titleformat{\subsubsection}{\normalsize\bfseries}{\thesubsubsection}{0.6em}{}
\titlespacing*{\section}{0pt}{12pt}{6pt}
\titlespacing*{\subsection}{0pt}{9pt}{4pt}

\newtcolorbox{defbox}[1]{
  breakable, enhanced, colback=fondonuevo, colframe=azulnuevo,
  boxrule=0.5pt, left=6pt, right=6pt, top=4pt, bottom=4pt,
  title={\bfseries #1}, coltitle=white, colbacktitle=azulnuevo, fonttitle=\small}
\newtcolorbox{yavisto}{
  breakable, enhanced, colback=fondorep, colframe=grisrep,
  boxrule=0.4pt, left=6pt, right=6pt, top=4pt, bottom=4pt, fontupper=\small}
\newtcolorbox{aviso}[1]{
  breakable, enhanced, colback=fondoaviso, colframe=ambar,
  boxrule=0.5pt, left=6pt, right=6pt, top=4pt, bottom=4pt,
  title={\bfseries #1}, coltitle=white, colbacktitle=ambar, fonttitle=\small}
\newtcolorbox{ejemplo}[1]{
  breakable, enhanced, colback=fondoej, colframe=verde,
  boxrule=0.5pt, left=6pt, right=6pt, top=4pt, bottom=4pt,
  title={\bfseries #1}, coltitle=white, colbacktitle=verde, fonttitle=\small}
\newtcolorbox{enresumen}{
  breakable, enhanced, colback=white, colframe=azulnuevo,
  boxrule=0.4pt, leftrule=3pt, left=6pt, right=6pt, top=4pt, bottom=4pt,
  title={\bfseries En resumen}, coltitle=azulnuevo, colbacktitle=white,
  fonttitle=\small, attach title to upper={\ }, fontupper=\small}

\setlist[itemize]{itemsep=1pt, topsep=2pt, parsep=0pt}
\setlist[enumerate]{itemsep=1pt, topsep=2pt, parsep=0pt}
\captionsetup{font=small, labelfont=bf}
\setlength{\parskip}{3pt}
\setlength{\parindent}{0pt}

\tikzset{
  caja/.style={draw=azulnuevo, fill=fondonuevo, rounded corners=2pt, align=center, font=\small, inner sep=4pt},
  flecha/.style={-{Stealth[length=2.2mm]}, thick, draw=azulnuevo},
  tarea/.style={draw=black!80, thick, minimum width=17mm, minimum height=8mm, fill=white, font=\small},
  crit/.style={tarea, draw=critrojo, fill=critrojo!15, very thick},
  temp/.style={text=azultemp, font=\scriptsize\bfseries},
  tard/.style={text=rojotard, font=\scriptsize\bfseries},
  dep/.style={-{Stealth[length=2.2mm]}, thick, draw=azulnuevo!75},
  depc/.style={-{Stealth[length=2.2mm]}, very thick, draw=critrojo}
}
\newcommand{\tnode}[9][tarea]{\node[#1,
  label={[temp,anchor=south west,inner sep=1pt]north west:#6},
  label={[temp,anchor=south east,inner sep=1pt]north east:#7},
  label={[tard,anchor=north west,inner sep=1pt]south west:#8},
  label={[tard,anchor=north east,inner sep=1pt]south east:#9}] (#2) at (#3) {\textbf{#4}\hspace{5mm}#5};}

\definecolor{fondoenun}{HTML}{F8FAFC}
\newtcolorbox{enunciado}[1]{
  breakable, enhanced, colback=fondoenun, colframe=black!55,
  boxrule=0.5pt, left=6pt, right=6pt, top=4pt, bottom=4pt,
  title={\bfseries #1}, coltitle=white, colbacktitle=black!65, fonttitle=\small}
\newtcolorbox{plantilla}[1]{
  breakable, enhanced, colback=fondorep, colframe=grisrep,
  boxrule=0.5pt, left=6pt, right=6pt, top=4pt, bottom=4pt,
  title={\bfseries Dado en el enunciado: #1}, coltitle=white, colbacktitle=grisrep, fonttitle=\small, fontupper=\small}
\newtcolorbox{resultado}[1]{
  breakable, enhanced, colback=fondoej, colframe=verde,
  boxrule=0.8pt, left=6pt, right=6pt, top=4pt, bottom=4pt,
  title={\bfseries #1}, coltitle=white, colbacktitle=verde, fonttitle=\small}
\newcommand{\hueco}[1][1.6cm]{\rule{0pt}{#1}}
\title{\vspace{-1.6cm}\textbf{Guía de ejercicios para el examen}\\[4pt]
\large Ejercicios del estilo de los TPs, con las fórmulas y plantillas que da el enunciado}
\author{\normalsize Administración de Proyectos de Software (7502) --- UNS}
\date{\normalsize Las respuestas están en \texttt{respuestas-ejercicios.pdf}}

\begin{document}
\maketitle
\thispagestyle{empty}

\begin{yavisto}
\textbf{Cómo usar esta guía.} Los ejercicios imitan a los de los TPs (el docente dijo que los del examen van a
ser ``muy similares''). Se priorizaron los tipos que la cátedra resolvió en clase: la red con los cuatro tipos
de dependencias (TP3, ejercicio 3), los Puntos de Función (TP4, sistema S1), la matriz persona/perfil (TP5,
ejercicio 2) y la tabla de riesgos con RSGR (TP1, caso de salud). Como en el parcial, cada enunciado trae en una
caja gris la \textbf{fórmula o plantilla} que hace falta: no hay que memorizarlas, hay que saber usarlas. No se
pide redactar un POS ni un COS, pero sí sacar requerimientos y riesgos de un enunciado largo (bloque B).
Resolvé primero y después compará con \texttt{respuestas-ejercicios.pdf}. La teoría de cada bloque está en el
resumen final (\texttt{examen/resumen/resumen-final.pdf}).
\end{yavisto}

\tableofcontents
\newpage

% =====================================================================
\section{Bloque A --- Scrum: historias, velocity y criterios de aceptación}
% =====================================================================
\begin{enunciado}{Ejercicio A1 (estilo TP0)}
Un club deportivo mantiene el siguiente Product Backlog para su sistema de reservas. ``Valor'' es el
\emph{business value} definido por el Product Owner y ``Esfuerzo'' son los \emph{story points} estimados por el
equipo.
\begin{center}\small
\begin{tabular}{@{}clcc@{}}
\toprule
\# & Ítem & Valor & Esfuerzo\\
\midrule
1 & Como socio, quiero reservar una cancha de pádel para asegurarme un horario & 8 & 5\\
2 & Como socio, quiero cancelar mi reserva para liberar el turno si no puedo ir & 5 & 2\\
3 & Migrar la base de datos a PostgreSQL 15 & 3 & 8\\
4 & Como administrativo, quiero ver la ocupación semanal de las canchas & 6 & 8\\
5 & Corregir el error de horarios duplicados en el calendario & 4 & 3\\
6 & Como socio, quiero recibir un recordatorio el día anterior a mi reserva & 3 & 3\\
7 & Configurar copias de seguridad automáticas & 2 & 5\\
8 & Como socio, quiero pagar la reserva con tarjeta & 7 & 13\\
9 & Como profesor, quiero publicar los horarios de mis clases & 2 & 5\\
\bottomrule
\end{tabular}
\end{center}
\begin{enumerate}[label=\alph*)]
  \item Describa brevemente el sistema que se está desarrollando.
  \item Indique qué ítems son historias de usuario y cuáles son ítems propios del equipo de desarrollo.
  Justifique.
  \item La velocity del equipo es 21. Indique qué ítems entran en el próximo sprint y en qué orden se eligen.
  \item Escriba los criterios de aceptación de las historias \#2 y \#6.
  \item El PO quiere agregar el ítem \#8 a mitad del sprint porque un sponsor lo pidió. ¿Qué le responde?
\end{enumerate}
\end{enunciado}
\begin{plantilla}{criterio de selección y formato de historia}
Ratio $=$ valor/esfuerzo; se toman los ítems de mayor a menor ratio mientras la suma de esfuerzos no supere la
velocity. Historia de usuario: ``Como [rol], quiero [objetivo], para [beneficio]''. Criterios de aceptación:
qué puede hacer el usuario; reglas ``si\ldots entonces\ldots''; confirmación visible al final.
\end{plantilla}

% =====================================================================
\section{Bloque B --- Requerimientos y riesgos a partir de un enunciado largo}
% =====================================================================
\begin{enunciado}{Ejercicio B1 (estilo TP1 + TP2): turnos en centros de salud municipales}
\small
La Secretaría de Salud de un municipio de 180.000 habitantes administra \textbf{12 centros de atención
primaria}, tres de ellos en zonas rurales con conectividad a Internet irregular. Hoy los turnos se piden por
teléfono o en persona: en las horas pico la espera telefónica supera los \textbf{40 minutos} y el
\textbf{25\% de los turnos otorgados no se usa} porque los pacientes se olvidan o no avisan que no van a ir. Los
centros atienden unas \textbf{85.000 consultas por año}.

El Secretario de Salud pidió ``un sistema de turnos como el que usa la ciudad vecina, porque ese anda bien''.
Después de dos entrevistas con el equipo de la Secretaría se sabe que:
\begin{itemize}
  \item Los pacientes deben poder pedir, consultar y cancelar turnos por una página web y por una app móvil,
  y los administrativos de cada centro deben poder darlos también en mostrador para quienes no usan Internet.
  \item Cada profesional tiene una agenda semanal que define el director de cada centro; los feriados y
  licencias bloquean la agenda.
  \item El sistema debe enviar un recordatorio por SMS 48 horas antes del turno, y si el paciente cancela, el
  turno se ofrece automáticamente al primero de una lista de espera.
  \item Solo pueden sacar turno los pacientes empadronados en el \textbf{padrón provincial de beneficiarios},
  que mantiene el Ministerio de Salud provincial; el sistema solo lo consulta.
  \item La Secretaría quiere un reporte mensual con la cantidad de turnos dados, usados y perdidos por centro y
  por especialidad, y el porcentaje de ausentismo.
  \item ``El sistema tiene que ser rápido y fácil de usar'', porque muchos pacientes son adultos mayores.
  \item Los datos de los pacientes son datos de salud y deben protegerse según la ley de protección de datos
  personales.
  \item El sistema debe estar funcionando antes del invierno (en \textbf{8 meses}), cuando empieza la campaña de
  vacunación. Se financia con un programa provincial cuyo presupuesto es fijo y vence al terminar el año.
  \item El área de informática municipal tiene \textbf{3 personas}; ninguna desarrolló antes una app móvil, y una
  de ellas tiene una licencia por maternidad prevista para dentro de 4 meses.
  \item Los administrativos de los centros temen que el sistema les agregue trabajo. Dentro de 10 meses hay
  elecciones municipales.
\end{itemize}
\normalsize
\begin{enumerate}[label=\alph*)]
  \item Liste los requerimientos funcionales del sistema (al menos 8). Para dos de ellos complete la ficha de
  requerimiento funcional.
  \item Identifique los requerimientos no funcionales. Elija uno que sea un atributo crítico y especifíquelo con
  la plantilla de atributos (proponga valores razonables y justifíquelos).
  \item En el enunciado hay al menos una \emph{solución disfrazada de requerimiento} y un \emph{requerimiento
  ambiguo}. Identifíquelos y explique cómo los trataría.
  \item Construya una lista de comprobación de elementos de riesgo (entre 10 y 15), indicando para cada uno la
  categoría de Pressman y, si corresponde, el conductor de riesgo de Wysocki.
  \item Elabore la tabla de riesgos justificando probabilidad e impacto, ordénela, defina una línea de corte y
  diga qué riesgos requieren mayor atención.
  \item Desarrolle el plan RSGR para los tres riesgos más prioritarios.
\end{enumerate}
\end{enunciado}

\begin{plantilla}{ficha de requerimiento funcional}
\begin{tabularx}{\linewidth}{@{}l X@{}}
req.\ id & identificador único (F1, F2\ldots)\\
categoría & funcional, de interfaz, de datos, de seguridad\ldots\\
descripción & qué debe hacer el sistema, quién lo hace\\
términos & glosario de los términos del dominio usados\\
justificación & por qué existe\\
prioridad & alta / media / baja\\
dependencias & otros requerimientos de los que depende\\
documentos & documentos de soporte\\
factibilidad & análisis preliminar de si es realizable\\
verificación & cómo se comprueba que se cumplió\\
\end{tabularx}
\end{plantilla}
\begin{plantilla}{especificación de un atributo (requerimiento no funcional)}
\begin{tabularx}{\linewidth}{@{}l X@{}}
Nombre & el atributo\\
Escala & unidad de medida: qué se mide\\
Test & procedimiento concreto: cómo se mide\\
Peor caso & peor nivel aceptable; por debajo, el proyecto fracasa en ese atributo\\
Nivel planificado & nivel que el proyecto se compromete a alcanzar\\
Mejor caso & mejor nivel alcanzable (solo informativo)\\
Nivel actual & nivel en el sistema o proceso existente\\
\end{tabularx}
\end{plantilla}
\begin{plantilla}{riesgos}
\textbf{Categorías de la lista de comprobación (Pressman)}: tamaño del producto; impacto en el negocio;
características del cliente; definición del proceso; entorno de desarrollo; tecnología a construir; tamaño y
experiencia de la plantilla.

\textbf{Conductores de riesgo (Wysocki)}: planificación demasiado ambiciosa; subestimación del presupuesto;
obligaciones de contrato mal interpretadas; estimación de tamaño inadecuada; hardware nuevo o poco familiar;
personal calificado no adecuado; plan de desarrollo inadecuado; métodos pobres de ingeniería; falta de soporte
político o de necesidad; análisis de riesgos inadecuado; expectativas poco realistas; procesos o tecnologías
poco familiares; modelo de desarrollo inadecuado; requerimientos definidos pobremente; estructura organizacional
inadecuada.

\textbf{Tabla de riesgos}:\quad \begin{tabular}{|l|l|l|l|}\hline Elemento de riesgo & Categoría & Probabilidad & Impacto\\\hline\end{tabular}

Impacto: 1 catastrófico, 2 crítico, 3 marginal, 4 despreciable. Ordenar por probabilidad y luego por impacto.

\textbf{Plan RSGR} (por riesgo): \emph{Reducción} (evitar y mitigar, antes) / \emph{Supervisión} (controlar que
se ejecuta la reducción y detectar señales, durante) / \emph{Gestión} (qué hacer si ocurrió).
\end{plantilla}

\begin{enunciado}{Ejercicio B2: plataforma para una cooperativa frutihortícola}
\small
Una cooperativa que agrupa a \textbf{60 productores} frutihortícolas quiere vender directamente a comercios y a
particulares de la región. Hoy los pedidos llegan por WhatsApp y teléfono, los anota una administrativa en un
cuaderno, y cada fin de mes se calcula a mano cuánto le corresponde a cada productor: el cálculo tarda una
semana y genera reclamos frecuentes.

La cooperativa pide ``una planilla compartida en Google para que los productores carguen su stock''. En las
reuniones se aclaró que necesitan:
\begin{itemize}
  \item que cada productor publique semanalmente la mercadería disponible (producto, cantidad, precio, fecha de
  cosecha), y que los productos se den de baja solos cuando superan los días de frescura de su tipo;
  \item que comercios y particulares hagan pedidos desde una web, eligiendo día y franja de entrega;
  \item cobrar con tarjeta por medio de una pasarela de pagos externa y emitir la factura electrónica a través
  del servicio de facturación del organismo recaudador;
  \item una app para los 4 repartidores, con la hoja de ruta del día y la confirmación de cada entrega;
  \item la liquidación mensual a cada productor (ventas menos la comisión de la cooperativa) y un reporte de
  ventas por producto y por productor;
  \item soportar la temporada alta de diciembre, en la que los pedidos se triplican.
\end{itemize}
El proyecto se financia con un subsidio que exige rendir gastos por hitos. Lo desarrollará una consultora
externa de 5 personas que nunca trabajó con cooperativas. Varios productores mayores nunca usaron una
aplicación y la presidenta de la cooperativa, impulsora del proyecto, termina su mandato en 6 meses.
\normalsize
\begin{enumerate}[label=\alph*)]
  \item Liste los requerimientos funcionales y no funcionales.
  \item ¿Qué hay detrás del pedido de la planilla de Google? ¿Cuál es el requerimiento real?
  \item Identifique entre 10 y 15 riesgos con su categoría, arme la tabla, ordénela, trace la línea de corte y
  desarrolle el RSGR de los dos primeros.
\end{enumerate}
\end{enunciado}

\begin{enunciado}{Ejercicio B3: funcional, no funcional o solución}
Clasifique cada enunciado como requerimiento \textbf{funcional} (F), \textbf{no funcional} (NF),
\textbf{solución} (S) o una \textbf{mezcla} que hay que separar. Cuando sea una solución, indique qué
pregunta haría y cuál podría ser el requerimiento real.
\begin{enumerate}
  \item El sistema permite al vendedor registrar una venta.
  \item El 95\% de las consultas se responde en menos de 2 segundos con 300 usuarios concurrentes.
  \item Hay que usar una base de datos Oracle.
  \item El sistema emite un reporte mensual de ventas por sucursal con totales.
  \item Queremos un chatbot para reducir las llamadas al centro de atención.
  \item El sistema debe estar disponible el 99,5\% del tiempo en horario comercial.
  \item Generar un manual de usuario fácil de entender.
  \item El sistema debe correr en los servidores Linux que ya tiene la empresa.
  \item El usuario puede recuperar su contraseña por correo electrónico.
  \item Un cajero nuevo debe poder cobrar sin ayuda después de 2 horas de capacitación.
\end{enumerate}
\end{enunciado}

% =====================================================================
\section{Bloque C --- Comparación de soluciones (estilo TP2, ejercicio 5)}
% =====================================================================
\begin{enunciado}{Ejercicio C1}
Para un problema se definieron cuatro atributos de calidad (A, B, C y D). Tres equipos propusieron soluciones,
y el análisis de impacto (en \% del nivel planificado de cada atributo que aporta cada solución) es:
\begin{center}
\begin{tabular}{c|ccc}
\toprule
 & S1 & S2 & S3\\
\midrule
A & 60 & 20 & 30\\
B & 15 & 50 & 30\\
C & 10 & 40 & 35\\
D & 35 & 20 & 25\\
\bottomrule
\end{tabular}
\end{center}
\begin{enumerate}[label=\alph*)]
  \item ¿Qué solución elegiría según su impacto en los atributos? ¿Por qué?
  \item ¿Hay algún atributo determinante para esa elección? Justifique.
  \item Los costos son: S1 el 30\% del presupuesto máximo, S2 el 50\% y S3 el 40\%. ¿Cuál es ahora la preferida?
  \item Se establece que el peor caso del atributo C exige que la solución aporte al menos el 15\% de su nivel
  planificado. ¿Cambia la decisión?
  \item Usted lidera el equipo de S2. ¿Qué podría hacer, legítimamente, para que su solución sea la preferida?
\end{enumerate}
\end{enunciado}
\begin{plantilla}{criterios de comparación}
Se suman algebraicamente los impactos: en atributos de calidad, más es mejor; en recursos, menos es mejor.
Razón de impactos $= \sum$ impactos de calidad $/\ \sum$ impactos de recursos (mayor es mejor). Paso 3 del
proceso de búsqueda de soluciones: eliminar las soluciones que producen el peor caso en algún atributo.
\end{plantilla}

% =====================================================================
\section{Bloque D --- Estimación}
% =====================================================================
\begin{plantilla}{fórmulas de estimación}
duración $=$ esfuerzo / porcentaje de dedicación; \quad 3 puntos: $E = (O + 4M + P)/6$.
Delphi: estiman en forma individual; justifican los que caen en los cuartiles exteriores; la estimación es el
promedio de la tercera pasada. Delphi de banda ancha: cada uno da $O$, $M$ y $P$; se eliminan los extremos de cada
columna, se promedia y se aplica la fórmula de 3 puntos.
\end{plantilla}
\begin{enunciado}{Ejercicio D1: Delphi}
Ocho estimadores dan, en la primera pasada, las siguientes duraciones (días) para una tarea: 12, 15, 9, 14, 22,
13, 16, 30. (a) ¿Quiénes deben justificar su estimación y por qué la técnica lo pide? (b) En la tercera pasada
las estimaciones son 14, 15, 15, 16, 17, 15, 16, 18. ¿Cuál es la estimación del grupo? (c) ¿Qué ventaja tiene
Delphi frente a que el LP estime solo, y en qué se parece y diferencia de Planning Poker?
\end{enunciado}
\begin{enunciado}{Ejercicio D2: 3 puntos}
Para cuatro tareas se obtuvieron estas estimaciones (días):
\begin{center}\small
\begin{tabular}{lccc}
\toprule
Tarea & Optimista & Más probable & Pesimista\\
\midrule
T1 & 2 & 4 & 12\\
T2 & 5 & 8 & 17\\
T3 & 1 & 3 & 5\\
T4 & 10 & 12 & 26\\
\bottomrule
\end{tabular}
\end{center}
(a) Calcule la estimación de 3 puntos de cada tarea. (b) ¿Por qué en T1 la estimación es mayor que el valor más
probable? (c) ¿Cuál es la tarea más incierta?
\end{enunciado}
\begin{enunciado}{Ejercicio D3: Delphi de banda ancha}
Cinco estimadores dan (días):
\begin{center}\small
\begin{tabular}{lccc}
\toprule
Estimador & $O$ & $M$ & $P$\\
\midrule
1 & 4 & 6 & 10\\
2 & 3 & 7 & 12\\
3 & 5 & 6 & 9\\
4 & 2 & 5 & 20\\
5 & 4 & 9 & 11\\
\bottomrule
\end{tabular}
\end{center}
Calcule la estimación descartando en cada columna el valor mínimo y el máximo.
\end{enunciado}
\begin{enunciado}{Ejercicio D4: duración, esfuerzo y crash point}
Una tarea requiere 96 horas-persona; la jornada es de 8 horas. (a) ¿Cuánto dura con una persona al 100\%? (b) ¿Y
con dos personas al 60\% cada una? (c) Si debe terminar en 6 días, ¿cuántas personas a tiempo completo hacen
falta? (d) ¿Por qué no se puede hacer en 1 día con 12 personas? (e) ¿De qué depende el costo de la tarea y de
qué su fecha de fin?
\end{enunciado}

% =====================================================================
\section{Bloque E --- Planificación: redes, planificaciones, márgenes y camino crítico}
% =====================================================================
\begin{plantilla}{convención y fórmulas de la cátedra}
Días numerados desde 1. CTe $= \max_{\text{pred}}$(FTe) $+ 1$ (1 si no tiene predecesoras); FTe $=$ CTe $+ d - 1$;
FTa $= \min_{\text{suc}}$(CTa) $- 1$ (fin del proyecto si no tiene sucesoras); CTa $=$ FTa $- d + 1$; Ventana
$V =$ FTa $-$ CTe $+ 1$; Margen total MT $= V - d$; Margen libre ML $= \min_{\text{suc}}$(CTe) $-$ FTe $- 1$
(fin del proyecto $+ 1$ si no tiene sucesoras).

Otras dependencias (predecesora $q$, sucesora $p$): \textbf{CC}: CTe($p$) $\ge$ CTe($q$), holgura CTe($p$) $-$
CTe($q$); \textbf{FF}: FTe($p$) $\ge$ FTe($q$), holgura FTe($p$) $-$ FTe($q$); \textbf{CF}: FTe($p$) $\ge$
CTe($q$), holgura FTe($p$) $-$ CTe($q$). Si una tarea solo tiene predecesoras FF o CF: CTe $= 1$ y FTe $=$
máx(CTe $+ d - 1$, lo que impone la dependencia).

Nodo: \quad\begin{tikzpicture}[baseline=-1mm]\tnode{X}{0,0}{X}{d}{CTe}{FTe}{CTa}{FTa}\end{tikzpicture}
\end{plantilla}

\begin{enunciado}{Ejercicio E1 (estilo TP3, ejercicio 1)}
Construya la red PDM, calcule las planificaciones temprana y tardía, la ventana y los márgenes de cada tarea, y
determine el camino crítico.
\begin{center}\small
\begin{tabular}{ccl@{\qquad}ccl}
\toprule
Tarea & Dur. & Dependencias & Tarea & Dur. & Dependencias\\
\midrule
A & 5 & -- & E & 2 & B [FC]\\
B & 3 & -- & F & 7 & C [FC], D [FC]\\
C & 4 & A [FC] & G & 3 & E [FC]\\
D & 6 & A [FC], B [FC] & H & 4 & F [FC], G [FC]\\
\bottomrule
\end{tabular}
\end{center}
\end{enunciado}
\begin{enunciado}{Ejercicio E2 (estilo TP3, ejercicio 2: dependencias CC)}
Ídem para:
\begin{center}\small
\begin{tabular}{ccl@{\qquad}ccl}
\toprule
Tarea & Dur. & Dependencias & Tarea & Dur. & Dependencias\\
\midrule
A & 2 & -- & E & 3 & D [CC]\\
B & 4 & -- & F & 8 & C [FC]\\
C & 6 & A [FC], B [FC] & G & 4 & E [FC]\\
D & 5 & B [CC] & H & 2 & F [FC], G [FC]\\
\bottomrule
\end{tabular}
\end{center}
\end{enunciado}
\begin{enunciado}{Ejercicio E3 (estilo TP3, ejercicio 3, el que resolvió la cátedra: los cuatro tipos)}
(a) Elabore el diagrama de red PDM. (b) Obtenga el camino crítico. (c) Calcule las planificaciones temprana y
tardía. (d) Calcule el margen libre, el margen total y la ventana de cada tarea.
\begin{center}\small
\begin{tabular}{ccl@{\qquad}ccl}
\toprule
Tarea & Dur. & Dependencias & Tarea & Dur. & Dependencias\\
\midrule
A & 4 & -- & G & 7 & E [FC]\\
B & 6 & -- & H & 2 & G [FF]\\
C & 5 & A [FC] & I & 5 & E [FC]\\
D & 8 & B [FC] & J & 6 & I [CF]\\
E & 3 & C [FC], D [FC] & K & 9 & F [FC], I [FC]\\
F & 4 & D [CC] & L & 4 & G [FC]\\
\bottomrule
\end{tabular}
\end{center}
\end{enunciado}
\begin{enunciado}{Ejercicio E4: análisis de atrasos y compresión}
Sobre la red del ejercicio E1: (a) ¿qué pasa con el proyecto si D se atrasa 2 días?, ¿y si G se atrasa 8 días?, ¿y
si C se atrasa 3 días? (b) Si E se atrasa 1 día, ¿a qué otra tarea afecta y por qué? (c) El cliente pide terminar
en 20 días. ¿Sobre qué tareas actuaría, con qué técnicas y qué riesgos tiene?
\end{enunciado}

% =====================================================================
\section{Bloque F --- Puntos de Función y COCOMO}
% =====================================================================
\begin{plantilla}{Puntos de Función}
\begin{minipage}[t]{0.45\linewidth}
\begin{tabular}{lccc}
\toprule
Comp. & Simple & Prom. & Compl.\\
\midrule
EE & 3 & 4 & 6\\
SE & 4 & 5 & 7\\
CE & 3 & 4 & 6\\
AI & 7 & 10 & 15\\
AE & 5 & 7 & 10\\
\bottomrule
\end{tabular}
\end{minipage}\hfill
\begin{minipage}[t]{0.52\linewidth}
PFNA $= \sum$ cantidad $\times$ peso. Cada factor F1--F14 se califica de 0 (irrelevante) a 5 (esencial); TI es su
suma. FCT $= 0{,}65 + 0{,}01 \times$ TI. PFA $=$ PFNA $\times$ FCT.
\end{minipage}
\end{plantilla}
\begin{enunciado}{Ejercicio F1 (estilo TP4)}
Para cada sistema calcule (a) los puntos de función no ajustados y (b) el factor de complejidad técnica y los
puntos de función ajustados.
\begin{center}\small
\begin{tabular}{@{}l l l l@{}}
\toprule
 & S4 & S5 & S6\\
\midrule
Entradas externas & 6 simples & 5 promedio y 2 complejas & 8 simples\\
Salidas externas & 5 promedio & 4 complejas & 3 simples y 3 complejas\\
Consultas externas & 3 simples y 2 promedio & 6 promedio & 5 simples\\
Archivos internos & 4 promedio & 2 promedio y 3 complejos & 2 simples\\
Archivos externos & 1 simple & 2 promedio & 3 complejos\\
\bottomrule
\end{tabular}

\smallskip
\begin{tabular}{@{}lllllll@{}}
\toprule
 & Irrelevante (0) & 1 & 2 & Promedio (3) & 4 & Esencial (5)\\
\midrule
S4 & F2, F13 & F10 & F1, F11 & F3--F6, F12 & F7, F8, F14 & F9\\
S5 & & F11, F13 & F12 & F1, F2, F10 & F3--F7 & F8, F9, F14\\
S6 & F1--F5 & & F6--F9 & & F10--F12 & F13, F14\\
\bottomrule
\end{tabular}
\end{center}
(c) ¿Por qué S6 queda con menos puntos ajustados que no ajustados? ¿Qué valor del total de influencias deja los
PF sin cambio?
\end{enunciado}
\begin{enunciado}{Ejercicio F2: contar desde una especificación}
Un sistema de biblioteca universitaria debe: dar de alta, modificar y dar de baja socios; registrar préstamos;
registrar devoluciones (que actualizan el préstamo); permitir consultar el catálogo por autor o título y ver la
ficha de un socio, mostrando los datos tal como están guardados; emitir un comprobante de multa con el monto
calculado según los días de atraso; y emitir un reporte mensual de préstamos por categoría con totales. El
sistema mantiene los datos de socios, de libros y ejemplares, y de préstamos. Además valida que el socio sea
alumno consultando el padrón de alumnos que mantiene el sistema académico de la universidad.

Las complejidades ya se determinaron: los ABM de socios, la consulta del catálogo, la ficha del socio, el
comprobante de multa, el archivo de socios, el de préstamos y el padrón son \emph{simples}; el registro de
préstamos, el de devoluciones, el reporte mensual y el archivo de libros y ejemplares son \emph{promedio}.
(a) Clasifique cada componente (EE, SE, CE, AI o AE) y justifique los casos dudosos. (b) Calcule los PFNA. (c) Con
TI $= 40$, calcule los PFA.
\end{enunciado}

\begin{plantilla}{COCOMO}
$E = a \cdot S^{\,b} \cdot F$ (personas-mes), $D = c \cdot E^{\,d}$ (meses), personas $= E/D$; $S$ en miles de
líneas. Básico: $F = 1$. Intermedio: $F =$ producto de los multiplicadores de los conductores de costo.
\begin{center}
\begin{tabular}{lcccc@{\qquad}l}
\toprule
Modo (básico) & $a$ & $b$ & $c$ & $d$ & Intermedio: $a$\\
\midrule
Orgánico & 2,4 & 1,05 & 2,5 & 0,38 & 3,2\\
Semi-embebido & 3,0 & 1,12 & 2,5 & 0,35 & 3,0\\
Embebido & 3,6 & 1,20 & 2,5 & 0,32 & 2,8\\
\bottomrule
\end{tabular}
\end{center}
Conversión: 128 líneas de C por punto de función. Puntos objeto: pantalla 1/2/3, reporte 2/5/8 (simple/media/alta),
componente 3GL 10; NPO $=$ PO $\times (100 - \%\text{reuso})/100$; PM $=$ NPO/PROD, con PROD $=$ 4, 7, 13, 25, 50
(muy baja a muy alta).
\end{plantilla}
\begin{enunciado}{Ejercicio F3: COCOMO básico}
Un sistema tiene 120 puntos de función ajustados y se programará en C. (a) Estime su tamaño en KLOC. (b) Calcule
esfuerzo, duración y personas promedio en modo orgánico y en modo embebido. (c) Compare los resultados: ¿qué
crece más, el esfuerzo o la duración?, ¿por qué? (d) ¿En qué modo clasificaría un sistema de liquidación de
sueldos y uno de control de semáforos de una ciudad?
\end{enunciado}
\begin{enunciado}{Ejercicio F4: COCOMO intermedio}
El mismo sistema, en modo orgánico, con estos conductores (el resto nominales): RELY alta (1,15), CPLX muy alta
(1,30), ACAP alta (0,86), LEXP baja (1,07), TOOL alta (0,91) y SCED muy baja (1,23). (a) Calcule el factor de
ajuste, el esfuerzo y la duración. (b) ¿Qué conductores aumentan el esfuerzo y cuáles lo reducen? ¿Por qué SCED
``muy baja'' lo aumenta?
\end{enunciado}
\begin{enunciado}{Ejercicio F5: puntos objeto (COCOMO II)}
Una aplicación tiene 8 pantallas simples, 4 medias y 2 altas; 3 reportes simples, 2 medios y 1 alto; y 3
componentes 3GL. Se estima un 25\% de reuso y una productividad nominal. Calcule los puntos objeto, los nuevos
puntos objeto y el esfuerzo. ¿Cuánto sería con productividad baja?
\end{enunciado}

% =====================================================================
\section{Bloque G --- Lanzamiento: matriz por perfil, asignación y cambios}
% =====================================================================
\begin{plantilla}{matriz por perfil}
(1) Matriz Persona $\times$ Perfil; (2) matriz Actividad $\times$ Perfil; (3) combinación: matriz Persona $\times$
Actividad (candidatos de cada actividad, se elige uno). Un recurso es crítico si es el único que cubre un perfil
necesario.
\end{plantilla}
\begin{enunciado}{Ejercicio G1 (estilo TP5, ejercicio 2, resuelto por la cátedra)}
Este es el equipo disponible:
\begin{center}\small
\begin{tabular}{llr}
\toprule
Nombre & Perfil & Costo por semana\\
\midrule
Ana Torres & Analista funcional & 9.000\\
Bruno Díaz & Desarrollador backend & 11.000\\
Carla Gómez & Desarrolladora backend & 10.500\\
Diego Paz & Desarrollador frontend & 10.000\\
Elena Ruiz & Tester & 7.500\\
Federico Lima & Tester & 8.000\\
Gabriela Sosa & Administradora de bases de datos y desarrolladora backend & 12.000\\
Hugo Vera & Diseñador UX y desarrollador frontend & 9.500\\
\bottomrule
\end{tabular}
\end{center}
El proyecto tiene siete actividades: A1 relevamiento (analista), A2 diseño de la base de datos (DBA), A3 API de
turnos (backend), A4 pantallas (frontend), A5 prototipos de experiencia de usuario (UX), A6 pruebas (tester) y A7
servicio de notificaciones (backend). A3 y A7 se hacen en paralelo, y A4 y A5 también.
(a) Construya la matriz Persona/Perfil. (b) ¿Qué recursos son críticos? ¿Por qué? (c) Construya la matriz
Actividad/Perfil y combine ambas para asignar una persona a cada actividad al menor costo posible. ¿Cuánto cuesta
el equipo por semana? (d) ¿Qué haría con los recursos críticos?
\end{enunciado}
\begin{enunciado}{Ejercicio G2 (estilo TP5, ejercicio 1): asignar personas y horas}
Un proyecto de 8 semanas lo hace este equipo: Laura (LP), Martín (programador/analista, aspirante a LP),
Sofía (programadora), Tomás (tester) y Rocío (referente del cliente, hasta 10 h por semana). Cada uno trabaja
como máximo 40 h por semana. Tareas:
\begin{center}\small
\begin{tabular}{llll}
\toprule
Tarea & Semanas & Esfuerzo & Perfil\\
\midrule
T1 Gestión del proyecto & 1--8 & 8 h por semana & LP\\
T2 Relevamiento & 1--2 & 40 h $+$ 10 h del cliente & analista\\
T3 Diseño & 2--3 & 50 h & analista o programador\\
T4 Módulo de inscripciones & 3--6 & 120 h & programador\\
T5 Módulo de reportes & 3--6 & 100 h & programador\\
T6 Pruebas & 6--7 & 60 h $+$ 10 h del cliente & tester\\
T7 Capacitación e implantación & 8 & 30 h $+$ 8 h del cliente & analista\\
\bottomrule
\end{tabular}
\end{center}
(a) Asigne personas a las tareas y arme una tabla con las horas de cada persona en cada tarea. (b) Verifique que
nadie supere su carga semanal. (c) Justifique la asignación de Martín.
\end{enunciado}
\begin{enunciado}{Ejercicio G3: pedido de cambio de alcance}
En la semana 5 de un proyecto de 12 semanas, el cliente pide que el sistema de reservas también permita pagar
online. El equipo estima 3 semanas de un desarrollador. (a) Describa cómo se procesa el pedido. (b) Proponga
tres alternativas e indique a cuál de los seis resultados posibles del análisis de impacto corresponde cada una.
(c) ¿Quién elige la alternativa?
\end{enunciado}
\begin{enunciado}{Ejercicio G4: organización del equipo}
(a) Un proyecto necesita 12 programadores. ¿Qué estructura usaría y por qué no la estructura I? (b) En una
empresa matricial, la gerente de línea de una desarrolladora le pide que vuelva antes a su grupo. ¿Qué problema
de la organización matricial muestra este caso? (c) Una consultora contrata a un especialista externo por 3
meses para una integración. ¿Qué componente del personal es y qué riesgos trae?
\end{enunciado}

% =====================================================================
\section{Bloque H --- Monitoreo y calidad}
% =====================================================================
\begin{plantilla}{fórmulas}
Valor ganado: $SV = EV - PV$; $CV = EV - AC$; $SPI = EV/PV$; $CPI = EV/AC$. Densidad de defectos $=$
defectos/KLOC. System spoilage $=$ tiempo de corrección post-entrega / tiempo total de desarrollo. Fog $= 0{,}4
\times$ (palabras/oraciones $+ 100 \times$ palabras de 3 o más sílabas/palabras). Grafo del equipo: densidad $=
e/[n(n-1)/2]$; nivel promedio $= 2e/n$. Riesgo de auditoría: $RA = RI \times RC \times RD$.
\end{plantilla}
\begin{enunciado}{Ejercicio H1: valor ganado}
Un proyecto de 10 semanas tiene un presupuesto (BAC) de \$200.000. Al final de la semana 4 el valor planificado
es \$70.000 y el costo real \$68.000. Las tareas T1 (\$20.000), T2 (\$15.000) y T3 (\$25.000) están terminadas; T4
(\$20.000) está al 50\% y T5 (\$10.000) no empezó. (a) Calcule EV contando solo las tareas terminadas. (b) Calcule
SV, CV, SPI y CPI e interprételos. (c) ¿Qué cambiaría si se
contara T4 al 50\%? ¿Por qué la cátedra sugiere contar solo las terminadas?
\end{enunciado}
\begin{enunciado}{Ejercicio H2: carta de tendencia de hitos}
Desvíos de los hitos en días (positivo $=$ antes de lo planeado, negativo $=$ después). ¿Qué patrón muestra cada
proyecto y qué hay que hacer? \quad P1: $+1, 0, -2, -3, -5, -6$. \quad P2: $-1, -2, -1, -2, -1, -1, -2$. \quad P3: $0,
-1, +8$.
\end{enunciado}
\begin{enunciado}{Ejercicio H3: densidad de defectos y spoilage}
Módulos: A (12 KLOC, 30 defectos), B (4 KLOC, 22 defectos), C (20 KLOC, 18 defectos). El sistema llevó 3.200 h de
desarrollo y en el primer año se usaron 480 h en corregir fallas reportadas. (a) Calcule la densidad de cada
módulo y la del sistema. ¿Qué módulo revisaría? (b) Calcule el spoilage. (c) ¿Por qué la densidad de C podría
engañar?
\end{enunciado}
\begin{enunciado}{Ejercicio H4: grafo del equipo}
Un equipo de 6 personas (LP, A, B, C, D, E) se comunica así: LP con A, B, C y D; A con B; B con C; C con D; D con E.
Calcule tamaño, nivel de comunicación del equipo, densidad, nivel individual de cada uno y nivel promedio.
Interprete.
\end{enunciado}
\begin{enunciado}{Ejercicio H5: índice Fog}
Calcule el índice Fog del texto: ``El equipo entrega el sistema en marzo. Los usuarios prueban cada pantalla
durante una semana. Después, el líder del proyecto aprueba la versión final.'' ¿Cómo se interpreta el resultado
en castellano?
\end{enunciado}
\begin{enunciado}{Ejercicio H6: riesgo de auditoría}
Un auditor acepta un riesgo de auditoría del 4\%. Estima un riesgo inherente de 0,8 y un riesgo de control de
0,25. (a) ¿Qué riesgo de detección necesita? (b) ¿Y si el riesgo de control fuera 0,5? ¿Qué hace el auditor?
\end{enunciado}
\begin{enunciado}{Ejercicio H7: clasificaciones}
(a) Tipo de mantenimiento: (1) se corrige un cálculo de intereses erróneo; (2) se adapta el sistema a una nueva
alícuota de IVA; (3) se agrega la exportación a Excel que pidieron los usuarios; (4) se reestructura un módulo
difícil de entender. (b) Tipo de control (preventivo, detectivo o correctivo): (1) quien carga un pago no puede
aprobarlo; (2) se comparan los totales del lote con la planilla en papel; (3) backup diario con procedimiento de
restauración; (4) revisión semanal del log de accesos. (c) Tipo de abuso: (1) un empleado consulta datos de un
famoso; (2) roban una notebook con la base de clientes; (3) un ransomware cifra los servidores; (4) un atacante
entra por una vulnerabilidad web. (d) ¿A qué factor de McCall afecta cada caso? (1) el sistema no se puede llevar
a otro sistema operativo; (2) cambiar un reporte lleva semanas; (3) usuarios sin permiso ven sueldos.
\end{enunciado}

% =====================================================================
\section{Bloque I --- Preguntas conceptuales (estilo ``conceptos importantes'')}
% =====================================================================
\begin{enumerate}\small
  \item ¿En qué consiste un proyecto? ¿Cuáles son sus parámetros o restricciones? ¿Cuáles son sus etapas?
  \item ¿Cómo se clasifican los proyectos de software? ¿Qué proceso es obligatorio en todos los tipos?
  \item ¿Qué es el POS? ¿Qué diferencia hay entre meta, objetivo y criterio de éxito? ¿Y entre FCE y criterio de éxito?
  \item ¿Qué son las condiciones de satisfacción y cuáles son sus etapas?
  \item ¿Cómo se clasifican los riesgos? ¿Qué es una lista de comprobación? ¿Qué son los conductores de riesgo?
  ¿En qué consiste un plan RSGR? ¿Con qué herramientas se desarrolla y aplica?
  \item Diferencie requerimiento funcional, no funcional y solución. ¿Qué es un atributo y cómo se especifica?
  \item Explique la diferencia entre actividad y tarea. ¿Qué es el crash point?
  \item ¿Sobre qué aspectos se estima? Describa los métodos de estimación de la duración.
  \item ¿Qué es el WBS, para qué sirve, cómo se construye y qué características deben tener sus tareas?
  \item ¿Qué tipos de diagramas de red hay? Ventajas y desventajas.
  \item Describa los cuatro tipos de dependencias y las restricciones que las originan.
  \item Defina planificación temprana y tardía, ventana, camino crítico, margen total y libre.
  \item ¿Qué describe la especificación? ¿Qué miden los PF y cómo se calculan los PFNA y los PFA?
  \item Diferencias entre COCOMO básico, intermedio y avanzado. ¿Qué vale $F$? ¿Cómo varían $a$, $b$, $c$, $d$ según
  el modo? Características y ejemplos de cada modo. Objetivos y diferencias de COCOMO II.
  \item Organizaciones del equipo de trabajo: tipos y diferencias. ¿Por qué planificar la organización?
  \item ¿En qué consiste la reunión de lanzamiento y para qué sirve? ¿Cómo se maneja un pedido de cambio de alcance?
  \item ¿Para qué medir atributos externos? ¿Qué utilidad y problemas tiene la densidad de defectos?
  \item Definiciones de calidad de Fenton; enfoques para monitorearla; modelo de McCall y cómo se analizan sus
  criterios; ventajas y desventajas de un modelo único.
  \item Bug vs.\ defecto; system spoilage; cómo medir usabilidad; tipos de mantenimiento; productividad y sus
  problemas; influencia de la estructura del equipo.
  \item Calidad según ISO; componentes de ISO 9001 (QMS, trazabilidad, desarrollo); calidad en ISO 9126.
  \item Niveles de CMM y cómo se pasa de uno a otro. ¿Qué es CMMI?
  \item Auditoría de SI: definición, objetivos, abuso computacional (tipos, consecuencias, ejemplos), controles,
  riesgo de auditoría y pasos.
\end{enumerate}

\end{document}
